\section{Scope and verification}
This blueprint covers the integration project, Sections 2--7 of
\href{https://arxiv.org/abs/0809.0473v2}{Integration in algebraically closed valued fields}.
Section 1 is introductory. The graph records selected mathematical dependencies,
not a complete expression-level dependency graph. Every green node has linked
Lean declarations at proof revision \texttt{862aa448b431}; definitions are
marked checked because their actual implementations compile. The prose is a
proof outline: the full linked types determine all hypotheses and conclusions.
\href{https://github.com/daviddhilbert/integration-in-ACVF/actions/runs/43851204054}{Matching proof CI passed}.

Lean is pinned to 4.35.0; Mathlib, special transformations and ACVF-QE are pinned
in the project manifest. The language has named RV inversion with $0^{-1}=0$,
one absorbing RV zero and the actual RV quotient kernel. The multiplicative
valuation order reverses the paper's value-group order. Canonical-model,
ACVF(0,0), genuine cardinal saturation, original-base smallness, VF field base,
VF/Gamma generation and named nontrivial Gamma values remain explicit where
required by the linked theorem types. No smallness of internally enlarged
real code sets is silently assumed.

RV[k] means the finite-to-one volume-map category, with arbitrary auxiliary RV
arity and possibly infinite underlying domain. Section 5.5 uses explicit RV
presentation charts. The final ordinary Section 6 APIs used here are in
\texttt{Section7.CompletedSection6}: they discharge the older universal
Section 5 input using checked proofs. Legacy conditional APIs remain available
but are not counted as additional closed theorems.

The measured degree-zero correction is essential: RV and VF objects retain
arbitrary nonzero definable RV forms, zero isomorphisms preserve the form
exactly, and products multiply both forms. Ordinary degree-zero categories
remain ordinary. Section 7 uses the corrected \texttt{Section7.HK} API,
not the legacy neutral-zero all-degree interface.
Lemma 7.5 has nonempty selected coordinates. Function interfaces have positive
VF dimension, and retained-coordinate projections are nonempty; total
integration is a separate operation. No pure Gamma/RES integration theorem
or new coarsening target is claimed by this project.

Approved wording and boundary corrections in 3.11, 3.24, 4.3 and the rectangle
observation are retained. For 3.17 ordinary-open-source least-radius existence
is proved; unconditional whole-field least-radius existence is not claimed.
See \repoFile{PAPER-CORRECTIONS.md}{paper corrections},
\repoFile{PROJECT-AUDIT.md}{project audit}, and
\repoFile{section\%207/COMPLETION.md}{Section 7 completion}.
The only permitted logical axioms in exported audits are
\texttt{propext}, \texttt{Classical.choice} and \texttt{Quot.sound}.



\section{Section 2: finite images and volume descent}


\begin{corollary}[2.15: finite Gamma images]

\label{s2-finite}

\lean{IntegrationInACVF.Section2.corollary_2_15}

\leanok

A map from a subset of Gamma to RV with full graph definable over the original real/Gamma base has finite image. Genuine saturation, an infinite cardinal, small original real base and ACVF(0,0) remain explicit.

\end{corollary}

\begin{proof}

\leanok

Apply the pinned finite-projection theorem to the full invariant graph.

\end{proof}

\begin{corollary}[2.34: positive-degree volume descent]

\label{s2-volume}

\lean{IntegrationInACVF.Section2.corollary_2_34}

\leanok

For an actual essential isomorphism of finite-volume RV objects of positive degree, a compatible measure-preserving lifted bijection implies the RV volume equation off an RV-small set and the Gamma volume equation at every essential label. The underlying isomorphism is the supplied one.

\end{corollary}

\begin{proof}

\uses{s2-finite}

\leanok

Construct the finite-extension and Gauss-graph input. Use Jacobian uniqueness and finite label fibres for the RV equation; restrict to fixed labels and descend constant-value polydiscs for the pointwise Gamma equation.

\end{proof}

\begin{corollary}[2.34: exact weighted degree zero]

\label{s2-zero}

\lean{IntegrationInACVF.Section2.WeightedZeroVolumeDescent.corollary_2_34_zero_weighted}

\leanok

In the corrected upstream measured degree-zero category, descend the lifted bijection while retaining its supplied ordinary label map and preserving the RV form exactly at every point. The ordinary degree-zero category is unchanged.

\end{corollary}

\begin{proof}

\uses{s2-volume}

\leanok

The checked upstream weighted-zero category supplies the actual pointwise measure equation; no positive-dimensional exceptional-set inference is used at degree zero.

\end{proof}

\section{Section 3: atomic sets and parametric balls}


\begin{definition}[3.1: complete types and atomic sets]

\label{s3-atomic}

\lean{IntegrationInACVF.Section3.GeneratesCompleteType, IntegrationInACVF.Section3.IsAtomic}

\leanok

A set generates a complete type over its parameter base if every definable test contains it or misses it. Atomicity adds definability of that set. These are separate notions, so later complete-type ambient sets need not be definable.

\end{definition}

\begin{lemma}[3.2: adjoining a ball name]

\label{s3-2}

\lean{IntegrationInACVF.Section3ACVF.AdjoinBall.lemma_3_2}

\leanok

Parametric-ball atomicity over the full real/Gamma/ball base is preserved when the actual input ball name is adjoined, under the stated nonalgebraicity of the output ball name.

\end{lemma}

\begin{proof}

\uses{s3-atomic}

\leanok

The linked implementation constructs the actual objects over the original interpreted base. Its full Lean type records the precise standing hypotheses; no stronger uniformity or parameter-smallness claim is implicit here.

\end{proof}

\begin{lemma}[3.3: Gamma parameters]

\label{s3-3}

\lean{IntegrationInACVF.Section3ACVF.FullGammaAtomicity.lemma_3_3}

\leanok

An actual definable atomic set remains atomic after adjoining any finite Gamma tuple, including mixed VF/RV tuple ambients. ACVF(0,0) and an originally definable nontrivial Gamma value are retained.

\end{lemma}

\begin{proof}

\uses{s3-2}

\leanok

The linked implementation constructs the actual objects over the original interpreted base. Its full Lean type records the precise standing hypotheses; no stronger uniformity or parameter-smallness claim is implicit here.

\end{proof}

\begin{lemma}[3.4: ball atomicity]

\label{s3-4}

\lean{IntegrationInACVF.Section3ACVF.CompleteTypeBallAtomicity.lemma_3_4}

\leanok

Every actual open or closed ball, with arbitrary initially unnamed radius, and the singleton and whole-field cases satisfy the complete-type ball-atomicity conclusion over the full interpreted base. Only the actual ball name is adjoined; the ambient complete type need not be definable.

\end{lemma}

\begin{proof}

\uses{s3-3}

\leanok

The linked implementation constructs the actual objects over the original interpreted base. Its full Lean type records the precise standing hypotheses; no stronger uniformity or parameter-smallness claim is implicit here.

\end{proof}

\begin{lemma}[3.6: primitive decomposition]

\label{s3-6}

\lean{IntegrationInACVF.Section3ACVF.PrimitiveDecomposition.lemma_3_6}

\leanok

An atomic unary VF set has an actual finite nonempty disjoint primitive decomposition with common radius, valuative centre and primitive kind. Individual components need not be definable over the original base.

\end{lemma}

\begin{proof}

\uses{s3-4}

\leanok

The linked implementation constructs the actual objects over the original interpreted base. Its full Lean type records the precise standing hypotheses; no stronger uniformity or parameter-smallness claim is implicit here.

\end{proof}

\begin{corollary}[3.7: coverage at a fixed radius]

\label{s3-7}

\lean{IntegrationInACVF.Section3ACVF.SameRadiusCoverage.corollary_3_7_gamma}

\leanok

The primitive decomposition gives whole-ball coverage at the radius of an inner ball, including the singleton and whole-field cases. An initially unnamed radius is adjoined using Gamma atomicity.

\end{corollary}

\begin{proof}

\uses{s3-3, s3-6}

\leanok

The linked implementation constructs the actual objects over the original interpreted base. Its full Lean type records the precise standing hypotheses; no stronger uniformity or parameter-smallness claim is implicit here.

\end{proof}

\begin{lemma}[3.8: primitive images]

\label{s3-8}

\lean{IntegrationInACVF.Section3ACVF.PrimitiveImages.lemma_3_8}

\leanok

Actual definable partial-domain VF maps preserve the recorded atomicity and one-component conditions on their images.

\end{lemma}

\begin{proof}

\uses{s3-6}

\leanok

The linked implementation constructs the actual objects over the original interpreted base. Its full Lean type records the precise standing hypotheses; no stronger uniformity or parameter-smallness claim is implicit here.

\end{proof}

\begin{lemma}[3.9 and 3.10: RV constancy]

\label{s3-9}

\lean{IntegrationInACVF.Section3ACVF.OpenBallConstancy.lemma_3_9, IntegrationInACVF.Section3ACVF.OpenBallConstancy.corollary_3_10}

\leanok

Definable RV-powerset families in arbitrary finite RV arity are constant on an atomic open ball; the companion corollary applies to arbitrary-radius open subballs of a definable atomic set.

\end{lemma}

\begin{proof}

\uses{s3-3, s3-6}

\leanok

The linked implementation constructs the actual objects over the original interpreted base. Its full Lean type records the precise standing hypotheses; no stronger uniformity or parameter-smallness claim is implicit here.

\end{proof}

\begin{lemma}[3.11: corrected separation]

\label{s3-11}

\lean{IntegrationInACVF.Section3ACVF.Separation.lemma_3_11}

\leanok

The separation conclusion holds with the accepted nonempty-set hypothesis. Both radii may be initially unnamed. The empty-set case is excluded as recorded in the paper corrections.

\end{lemma}

\begin{proof}

\uses{s3-6, s3-9}

\leanok

The linked implementation constructs the actual objects over the original interpreted base. Its full Lean type records the precise standing hypotheses; no stronger uniformity or parameter-smallness claim is implicit here.

\end{proof}

\begin{lemma}[3.13: images of atomic open balls]

\label{s3-13}

\lean{IntegrationInACVF.Section3ACVF.OpenImages.lemma_3_13}

\leanok

A definable VF image of an atomic open ball is either constant or an atomic open ball.

\end{lemma}

\begin{proof}

\uses{s3-8, s3-9}

\leanok

The linked implementation constructs the actual objects over the original interpreted base. Its full Lean type records the precise standing hypotheses; no stronger uniformity or parameter-smallness claim is implicit here.

\end{proof}

\begin{lemma}[3.14: closed-image obstruction]

\label{s3-14}

\lean{IntegrationInACVF.Section3ACVF.ClosedImageObstruction.lemma_3_14}

\leanok

An infinite one-component atomic target cannot be closed when its only ball imaginary is the source open-ball name. Original Gamma names are allowed, and the source radius need not be named.

\end{lemma}

\begin{proof}

\uses{s3-13}

\leanok

The linked implementation constructs the actual objects over the original interpreted base. Its full Lean type records the precise standing hypotheses; no stronger uniformity or parameter-smallness claim is implicit here.

\end{proof}

\begin{lemma}[3.15: complete-type open images]

\label{s3-15}

\lean{IntegrationInACVF.Section3ACVF.FullCompleteTypeOpenImages.lemma_3_15}

\leanok

The open-image conclusion extends to the full original real/Gamma/ball base, arbitrary open-ball radii and possibly undefinable complete-type ambients. The map may be defined on a larger definable domain.

\end{lemma}

\begin{proof}

\uses{s3-4, s3-14}

\leanok

The linked implementation constructs the actual objects over the original interpreted base. Its full Lean type records the precise standing hypotheses; no stronger uniformity or parameter-smallness claim is implicit here.

\end{proof}

\begin{lemma}[3.16: Gamma linearity]

\label{s3-16}

\lean{IntegrationInACVF.Section3ACVF.GammaLinearity.lemma_3_16}

\leanok

The actual map is Gamma-linear with coefficient the target/source radius ratio, using actual radius properties. Linearity is proved, rather than assumed.

\end{lemma}

\begin{proof}

\uses{s3-15}

\leanok

The linked implementation constructs the actual objects over the original interpreted base. Its full Lean type records the precise standing hypotheses; no stronger uniformity or parameter-smallness claim is implicit here.

\end{proof}

\begin{corollary}[3.17: least radius for ordinary open balls]

\label{s3-17}

\lean{IntegrationInACVF.Section3ACVF.LeastLinearRadiusExistence.corollary_3_17_ordinary}

\leanok

The ordinary-open-source result constructs the least linear radius and proves the exclusion under injectivity and the stated standing hypotheses. The whole-field source has a separate conditional interface: existence of a least radius for the identity on VF is not asserted.

\end{corollary}

\begin{proof}

\uses{s3-16}

\leanok

The linked implementation constructs the actual objects over the original interpreted base. Its full Lean type records the precise standing hypotheses; no stronger uniformity or parameter-smallness claim is implicit here.

\end{proof}

\begin{lemma}[3.18: RV linearity]

\label{s3-18}

\lean{IntegrationInACVF.Section3ACVF.RVLinearity.lemma_3_18}

\leanok

Actual RV linearity holds in the recorded standing scope. The proof uses an affine-error replacement for the printed endpoint inference.

\end{lemma}

\begin{proof}

\uses{s3-16}

\leanok

The linked implementation constructs the actual objects over the original interpreted base. Its full Lean type records the precise standing hypotheses; no stronger uniformity or parameter-smallness claim is implicit here.

\end{proof}

\begin{proposition}[3.19: open-to-open partition]

\label{s3-19}

\lean{IntegrationInACVF.Section3ACVF.FullOpenToOpenPartition.proposition_3_19}

\leanok

Construct a finite disjoint partition into nonempty original-base definable pieces on which the actual map is open-to-open. There is no cardinal bound on old imaginary names.

\end{proposition}

\begin{proof}

\uses{s3-15, s3-18}

\leanok

The linked implementation constructs the actual objects over the original interpreted base. Its full Lean type records the precise standing hypotheses; no stronger uniformity or parameter-smallness claim is implicit here.

\end{proof}

\begin{proposition}[3.21: coordinate partitions]

\label{s3-21}

\lean{IntegrationInACVF.Section3ACVF.CoordinatePartitions.proposition_3_21}

\leanok

Construct one uniform finite nonempty disjoint original-base partition for the selected VF coordinate pairs of actual mixed finite-tuple injections. Arbitrary finite index types cover all mixed arities.

\end{proposition}

\begin{proof}

\uses{s3-19}

\leanok

The linked implementation constructs the actual objects over the original interpreted base. Its full Lean type records the precise standing hypotheses; no stronger uniformity or parameter-smallness claim is implicit here.

\end{proof}

\begin{lemma}[3.22: components and holes]

\label{s3-22}

\lean{IntegrationInACVF.Section3ACVF.ThinFamilyConservation.lemma_3_22}

\leanok

An original-base definable open-to-open VF bijection preserves components and actual holes for finite disjoint thin families. This theorem needs ACVF and a genuine nontrivial value, but no saturation or characteristic-zero premise.

\end{lemma}

\begin{proof}

\uses{s3-19}

\leanok

The linked implementation constructs the actual objects over the original interpreted base. Its full Lean type records the precise standing hypotheses; no stronger uniformity or parameter-smallness claim is implicit here.

\end{proof}

\begin{proposition}[3.23: primitive distinction]

\label{s3-23}

\lean{IntegrationInACVF.Section3ACVF.AtomicPrimitiveDistinction.proposition_3_23}

\leanok

There is no definable bijection between distinct actual atomic primitive shapes: open ball, closed ball and thin annulus. ACVF(0,0) and an originally definable nontrivial value remain explicit.

\end{proposition}

\begin{proof}

\uses{s3-22}

\leanok

The linked implementation constructs the actual objects over the original interpreted base. Its full Lean type records the precise standing hypotheses; no stronger uniformity or parameter-smallness claim is implicit here.

\end{proof}

\begin{proposition}[3.24: corrected punctured-ball obstruction]

\label{s3-24}

\lean{IntegrationInACVF.Section3ACVF.PuncturedBallObstruction.proposition_3_24}

\leanok

The author-corrected obstruction holds for an ordinary closed parent with a single proper open hole. The open target includes VF. The full original interpreted base and displayed saturation, smallness and ACVF(0,0) hypotheses are retained.

\end{proposition}

\begin{proof}

\uses{s3-22, s3-23}

\leanok

The linked implementation constructs the actual objects over the original interpreted base. Its full Lean type records the precise standing hypotheses; no stronger uniformity or parameter-smallness claim is implicit here.

\end{proof}

\section{Section 4: preparation and two-cells}


\begin{lemma}[4.1: translation offsets]

\label{s4-1}

\lean{IntegrationInACVF.Section4.GammaOffsets.lemma_4_1}

\leanok

An actual special bijection has an original-Gamma-base offset graph and the exact inverse-fibre translation identity, including the zero case.

\end{lemma}

\begin{proof}

\uses{s3-9}

\leanok

Construct the preparation uniformly over the original base using the earlier atomic geometry and actual special transformations; the linked theorem retains its displayed model, base and cardinal hypotheses.

\end{proof}

\begin{lemma}[4.2: linear preparation]

\label{s4-2}

\lean{IntegrationInACVF.Section4.LinearPreparation.lemma_4_2}

\leanok

Construct an actual original-base special bijection to an RV pullback whose inverse VF pieces are Gamma-linear and open-to-open. Prepare the full radius/slope relation; no definable least-radius function is assumed.

\end{lemma}

\begin{proof}

\uses{s3-16, s3-19, s4-1}

\leanok

Construct the preparation uniformly over the original base using the earlier atomic geometry and actual special transformations; the linked theorem retains its displayed model, base and cardinal hypotheses.

\end{proof}

\begin{lemma}[4.3: corrected radius preparation]

\label{s4-3}

\lean{IntegrationInACVF.Section4.RadiusPreparation.lemma_4_3}

\leanok

One preparation works for every ordinary target radius, with at most one source radius on each actual inverse piece. The whole-field source is included. This is the accepted corrected statement.

\end{lemma}

\begin{proof}

\uses{s4-2}

\leanok

Construct the preparation uniformly over the original base using the earlier atomic geometry and actual special transformations; the linked theorem retains its displayed model, base and cardinal hypotheses.

\end{proof}

\begin{lemma}[4.4: inverse differences]

\label{s4-4}

\lean{IntegrationInACVF.Section4.InverseDifferencePreparation.lemma_4_4}

\leanok

Prepare the actual inverse-difference RV range to a singleton using translated RV fibres and an actual open-to-open positioning bijection. Individual fibre paradigms need not be definable.

\end{lemma}

\begin{proof}

\uses{s4-2, s4-3}

\leanok

Construct the preparation uniformly over the original base using the earlier atomic geometry and actual special transformations; the linked theorem retains its displayed model, base and cardinal hypotheses.

\end{proof}

\begin{lemma}[4.6: balanced slope]

\label{s4-6}

\lean{IntegrationInACVF.Section4.BalancedSlope.remark_4_6}

\leanok

For a nondegenerate balanced configuration the slope is the ratio of the two paradigm values, derived from the actual fibre identities and linearity.

\end{lemma}

\begin{proof}

\uses{s4-balanced, s3-16}

\leanok

Construct the preparation uniformly over the original base using the earlier atomic geometry and actual special transformations; the linked theorem retains its displayed model, base and cardinal hypotheses.

\end{proof}

\begin{lemma}[4.8: definable two-cell decomposition]

\label{s4-8}

\lean{IntegrationInACVF.Section4.lemma_4_8}

\leanok

Every original-base definable native VF-squared set admits an actual finite-arity Gamma RV map whose nonempty fibres are definable two-cells over the original Gamma base with the output RV tuple. The positioning graph and both paradigms are definable; no cell-existence or family-gluing conclusion is assumed.

\end{lemma}

\begin{proof}

\uses{s4-4, s4-cells, s3-21}

\leanok

Construct the preparation uniformly over the original base using the earlier atomic geometry and actual special transformations; the linked theorem retains its displayed model, base and cardinal hypotheses.

\end{proof}

\begin{definition}[4.5: balanced fibres]

\label{s4-balanced}

\lean{IntegrationInACVF.Section4.BalancedFibres.HasBalancedFibres}

\leanok

The two actual fibre identities specify a balanced relation, positioning map and RV paradigms. Degenerate paradigms are handled explicitly.

\end{definition}

\begin{definition}[4.7: definable two-cells]

\label{s4-cells}

\lean{IntegrationInACVF.Section4.Cells.DefinableTwoCell}

\leanok

A two-cell consists of actual one-cell projections and one of the recorded positioning alternatives, with the full original-Gamma definability of the relation, positioning map and paradigms.

\end{definition}

\begin{proof}

\uses{s4-balanced}

\leanok



\end{proof}

\begin{proposition}[Corrected rectangle observation]

\label{s4-rectangle}

\lean{IntegrationInACVF.Section4.CellRectangles.rectangle, IntegrationInACVF.Section4.CellRectangles.open_polydisc}

\leanok

A nonbalanced cell is a product of a one-cell and an open ball after positioning. It is an open polydisc when the first projection is open. The unrestricted open-polydisc wording is replaced by the accepted rectangle statement.

\end{proposition}

\begin{proof}

\uses{s4-cells}

\leanok

Identify both actual fibres; the point-times-open-ball counterexample explains why openness of the first projection is needed.

\end{proof}

\section{Section 5: finite-volume standard contractions}


\begin{lemma}[5.1: finite preparation]

\label{s5-1}

\lean{IntegrationInACVF.Section5.lemma_5_1_finite}

\leanok

Construct two actual Gamma special bijections to RV pullbacks, a contractible conjugate with its full Gamma graph, and finite volume maps for both prepared objects.

\end{lemma}

\begin{proof}

\uses{s4-2}

\leanok

Use the checked uniform preparation and finite contraction compiler. RV[k] means finite-to-one volume maps, with arbitrary auxiliary RV arity; the full theorem type retains saturation, original-base smallness and the required field, generation and named-value hypotheses.

\end{proof}

\begin{lemma}[5.2: comparison of unary preparations]

\label{s5-2}

\lean{IntegrationInACVF.Section5.lemma_5_2_finite}

\leanok

Construct both special preparations, mutual contractibility, the unique commuting key bijection and an actual RV[1] isomorphism with both Gamma graphs. Target finiteness follows from source finiteness.

\end{lemma}

\begin{proof}

\uses{s5-1}

\leanok

Use the checked uniform preparation and finite contraction compiler. RV[k] means finite-to-one volume maps, with arbitrary auxiliary RV arity; the full theorem type retains saturation, original-base smallness and the required field, generation and named-value hypotheses.

\end{proof}

\begin{lemma}[5.3: the finite-category square]

\label{s5-3}

\lean{IntegrationInACVF.Section5.lemma_5_3_finite}

\leanok

The prepared mutually contractible unary bijection induces the actual finite-category isomorphism and commuting key square.

\end{lemma}

\begin{proof}

\uses{s5-2}

\leanok

Use the checked uniform preparation and finite contraction compiler. RV[k] means finite-to-one volume maps, with arbitrary auxiliary RV arity; the full theorem type retains saturation, original-base smallness and the required field, generation and named-value hypotheses.

\end{proof}

\begin{lemma}[5.4: relative-special factorization]

\label{s5-4}

\lean{IntegrationInACVF.Section5.special_relativeSpecial_factorization}

\leanok

An actual special bijection factors through actual coordinate-fibre relative-special witnesses. Complementary VF coordinates become fibre parameters. Gamma bijections preserve finite RV fibres.

\end{lemma}

\begin{proof}

\uses{s5-1}

\leanok

Use the checked uniform preparation and finite contraction compiler. RV[k] means finite-to-one volume maps, with arbitrary auxiliary RV arity; the full theorem type retains saturation, original-base smallness and the required field, generation and named-value hypotheses.

\end{proof}

\begin{lemma}[5.5: presented standard contraction]

\label{s5-5}

\lean{IntegrationInACVF.Section5.lemma_5_5}

\leanok

For every VF order, including zero arity, construct a standard contraction with native and normalized RV[n] objects, a natural Gamma lift isomorphism and an actual global special presentation. Explicit original-Gamma RV charts connect the special-chain target to the actual lift; no unchanged-arity literal special witness is asserted.

\end{lemma}

\begin{proof}

\uses{s5-3, s5-4}

\leanok

Use the checked uniform preparation and finite contraction compiler. RV[k] means finite-to-one volume maps, with arbitrary auxiliary RV arity; the full theorem type retains saturation, original-base smallness and the required field, generation and named-value hypotheses.

\end{proof}

\begin{lemma}[5.6: relatively unary factors]

\label{s5-6}

\lean{IntegrationInACVF.Section5.lemma_5_6_finite}

\leanok

In positive VF arity construct a finite disjoint covering Gamma partition, actual relatively unary factors, a common real realization and finite RV fibres at every intermediate stage.

\end{lemma}

\begin{proof}

\uses{s5-4, s5-5}

\leanok

Use the checked uniform preparation and finite contraction compiler. RV[k] means finite-to-one volume maps, with arbitrary auxiliary RV arity; the full theorem type retains saturation, original-base smallness and the required field, generation and named-value hypotheses.

\end{proof}

\begin{lemma}[5.7: two-cell contractions]

\label{s5-7}

\lean{IntegrationInACVF.Section5.lemma_5_7_finite}

\leanok

Every definable two-cell has both actual standard contractions and an actual Gamma RV[2] isomorphism. Finite volume is derived without an additional source assumption.

\end{lemma}

\begin{proof}

\uses{s4-8, s5-5}

\leanok

Use the checked uniform preparation and finite contraction compiler. RV[k] means finite-to-one volume maps, with arbitrary auxiliary RV arity; the full theorem type retains saturation, original-base smallness and the required field, generation and named-value hypotheses.

\end{proof}

\begin{lemma}[5.8: two contraction orders]

\label{s5-8}

\lean{IntegrationInACVF.Section5.lemma_5_8_finite}

\leanok

Construct a VF-fixed positioned cell-label injection and both whole-source standard contractions with an actual Gamma RV[2] comparison and both finite volume maps.

\end{lemma}

\begin{proof}

\uses{s5-6, s5-7}

\leanok

Use the checked uniform preparation and finite contraction compiler. RV[k] means finite-to-one volume maps, with arbitrary auxiliary RV arity; the full theorem type retains saturation, original-base smallness and the required field, generation and named-value hypotheses.

\end{proof}

\section{Section 6: blowups and ordinary integration}


\begin{lemma}[6.3--6.6: finite common blowups]

\label{s6-confluence}

\lean{IntegrationInACVF.Section6.FiniteRV.lemma_6_3, IntegrationInACVF.Section6.FiniteRV.lemma_6_5, IntegrationInACVF.Section6.FiniteRV.corollary_6_6}

\leanok

Elementary blowups of isomorphic finite-volume RV objects are isomorphic. Iterated blowups admit common refinements in the actual finite category, yielding isomorphism-invariant common-blowup equivalence.

\end{lemma}

\begin{proof}

\uses{s5-5}

\leanok

Transport elementary blowups and commute successive steps, deriving finite volume for every intermediate object.

\end{proof}

\begin{lemma}[6.8: the blowup congruence]

\label{s6-congruence}

\lean{IntegrationInACVF.Section6.FiniteRV.lemma_6_8_add_class_iff, IntegrationInACVF.Section6.FiniteRV.lemma_6_8_mul_class_iff}

\leanok

On actual RV[k] representatives the additive and graded multiplicative congruences agree exactly with common-blowup equivalence. Cartesian products preserve this relation.

\end{lemma}

\begin{proof}

\uses{s6-confluence}

\leanok

Construct the congruences from actual finite partitions and elementary blowups.

\end{proof}

\begin{lemma}[6.10 and 6.12: special-chain lift diagrams]

\label{s6-special}

\lean{IntegrationInACVF.Section6.FiniteRV.corollary_6_10_gamma, IntegrationInACVF.Section6.lemma_6_12}

\leanok

Actual Gamma-certified special chains induce finite blowup diagrams. The elementary lift isomorphism is constructed over the original Gamma base, without a canonical-presentation assumption.

\end{lemma}

\begin{proof}

\uses{s6-confluence, s5-4}

\leanok

Use the actual focus and graph constructions, preserving finite intermediate objects and the genuine VF/Gamma-generated-base hypothesis where required.

\end{proof}

\begin{theorem}[Closing the Section 5 input]

\label{s6-closed-input}

\lean{IntegrationInACVF.Section7.CompletedSection6.lemma_6_17}

\leanok

The closed generated-base API proves the ordinary integration kernel equivalence from the completed Section 5 theorem, including its needed finite-VF-parameter instances. It carries no unproved universal Section 5.8 hypothesis.

\end{theorem}

\begin{proof}

\uses{s5-8, s6-congruence, s6-special}

\leanok

Descend the finite parameter bases through definable closure, prove contraction-order independence without deleting zero RV labels, and identify the kernel.

\end{proof}

\begin{proposition}[6.18: additive integration]

\label{s6-18}

\lean{IntegrationInACVF.Section7.CompletedSection6.proposition_6_18}

\leanok

Integration is an additive equivalence from actual finite Gamma VF[k] scissors to the finite RV[k] common-blowup quotient, with both inverse laws.

\end{proposition}

\begin{proof}

\uses{s6-closed-input}

\leanok

Use the closed kernel equivalence and actual lifts; prove representative independence by finite partition gluing.

\end{proof}

\begin{theorem}[6.19: global semiring equivalence]

\label{s6-19}

\lean{IntegrationInACVF.Section7.CompletedSection6.theorem_6_19}

\leanok

Global ordinary integration is a semiring equivalence with the graded RV blowup target. Addition is the original scissor addition and multiplication is the actual Cartesian product.

\end{theorem}

\begin{proof}

\uses{s6-18, s6-congruence}

\leanok

Construct the finite-stage congruence union and bilinear Cartesian product; integration preserves products and the unit.

\end{proof}

\begin{remark}[6.20: definable families]

\label{s6-20}

\lean{IntegrationInACVF.Section7.CompletedSection6.remark_6_20_representative_independence, IntegrationInACVF.Section7.CompletedSection6.remark_6_20_linearity}

\leanok

Actual parameter-dependent coefficient semirings over definable closure give representative-independent family integration. Integration is linear on the definable-function quotient, including zero, addition and scalar multiplication.

\end{remark}

\begin{proof}

\uses{s6-19}

\leanok

Construct literal total family lifts and compare arbitrary finite RV degrees.

\end{proof}

\begin{proposition}[6.21: coordinate Fubini]

\label{s6-21}

\lean{IntegrationInACVF.Section7.CompletedSection6.proposition_6_21}

\leanok

Coordinate pushforwards on actual definable functions have the literal fibre-integral formula, and iterated integration equals total integration.

\end{proposition}

\begin{proof}

\uses{s6-20}

\leanok

Descend the geometric coordinate maps through the function quotient using the proved family identities.

\end{proof}

\begin{theorem}[6.22: ordinary localization]

\label{s6-22}

\lean{IntegrationInACVF.Section7.CompletedSection6.theorem_6_22, IntegrationInACVF.Section7.CompletedSection6.theorem_6_22_injective}

\leanok

The groupified ordinary VF ring embeds into the localized ordinary RV ring, with the proved homogeneous normalization formula. Ordinary degree zero retains its original category.

\end{theorem}

\begin{proof}

\uses{s6-19}

\leanok

Identify the principal ordinary relation, construct the graded additive colimit and use injective padding. These are ordinary results, not the corrected measured principal-kernel claim of Section 7.

\end{proof}

\section{Section 7: volume forms and corrected class algebra}


\begin{definition}[Corrected RV volume objects]

\label{s7-zero-rv}

\lean{IntegrationInACVF.Section7.HK.ZeroRVVolumeObject, IntegrationInACVF.Section7.HK.RVVolumeObject}

\leanok

In every degree use an actual finite-volume RV object with a nonzero definable RV form. In measured degree zero retain that arbitrary form, with exact weight-preserving zero isomorphisms. Positive objects agree with the existing positive-degree objects; ordinary RV[0] is unchanged.

\end{definition}

\begin{proof}

\uses{s2-zero}

\leanok



\end{proof}

\begin{definition}[Corrected VF volume objects and zero integration]

\label{s7-zero-vf}

\lean{IntegrationInACVF.Section7.HK.VFVolumeObject, IntegrationInACVF.Section7.HK.zeroVolumeIntegrationEquiv}

\leanok

Measured VF degree zero likewise carries an arbitrary nonzero definable RV form. Its lift and integration equivalence preserve the actual form. Positive-dimensional objects retain the actual partial Jacobian measure condition.

\end{definition}

\begin{proof}

\uses{s7-zero-rv}

\leanok



\end{proof}

\begin{lemma}[Multiplication of both volume forms]

\label{s7-products}

\lean{IntegrationInACVF.Section7.HK.weighted_lift_zero_product_map_density}

\leanok

Products multiply both RV forms in every degree, including either factor of degree zero. Weighted zero lifting respects this actual density equation.

\end{lemma}

\begin{proof}

\uses{s7-zero-rv, s7-zero-vf}

\leanok

The zero-degree factor contributes its density; retaining this factor repairs the degree-one factorization failure of the legacy neutral-zero convention.

\end{proof}

\begin{lemma}[7.2: elementary volume blowup comparison]

\label{s7-2}

\lean{IntegrationInACVF.Section7.lemma_7_2}

\leanok

Elementary blowups of isomorphic actual positive-volume objects are volume-isomorphic, under the displayed canonical-model and base-context hypotheses.

\end{lemma}

\begin{proof}

\uses{s6-confluence}

\leanok

Construct and compare the elementary volume diagrams over the recorded base.

\end{proof}

\begin{lemma}[7.4: volume blowup congruence]

\label{s7-4}

\lean{IntegrationInACVF.Section7.lemma_7_4}

\leanok

The positive-volume semiring relation is identified with actual common volume blowups. Corrected zero-degree adapters retain exact weighted zero isomorphisms.

\end{lemma}

\begin{proof}

\uses{s7-2, s6-congruence}

\leanok

Use elementary comparisons, finite partitions and product compatibility.

\end{proof}

\begin{lemma}[7.5: specified selected contractions]

\label{s7-5}

\lean{IntegrationInACVF.Section7.lemma_7_5, IntegrationInACVF.Section7.lemma_7_5_ordinary}

\leanok

For the specified actual standard contractions, fibrewise common-volume-blowup equivalence over the common VF projection gives common blowups of the appended objects after the displayed base extension. The selected-coordinate set is nonempty; that condition is carried by the selected-family data.

\end{lemma}

\begin{proof}

\uses{s5-8, s7-4}

\leanok

Prepare the selected families uniformly and compare the supplied contractions. The empty selected set is outside the accepted statement.

\end{proof}

\begin{corollary}[7.6: contraction-order comparison]

\label{s7-6}

\lean{IntegrationInACVF.Section7.corollary_7_6}

\leanok

The actual positive-dimensional contractions have the recorded volume-blowup comparison, with the original selected volume labels and standing hypotheses retained.

\end{corollary}

\begin{proof}

\uses{s7-5}

\leanok

Apply the specified-contraction comparison to the actual family data.

\end{proof}

\begin{lemma}[7.7: volume coordinate exchange]

\label{s7-7}

\lean{IntegrationInACVF.Section7.lemma_7_7}

\leanok

The actual volume contraction exchange holds in the native and ordinary interfaces, supplying independence of the selected contraction order.

\end{lemma}

\begin{proof}

\uses{s7-6}

\leanok

Compare actual coordinate preparations and use common volume blowups.

\end{proof}

\begin{proposition}[7.8: positive volume lift kernel]

\label{s7-8}

\lean{IntegrationInACVF.Section7.proposition_7_8}

\leanok

For positive-degree measured objects, equality of the actual lifted classes is characterized by common volume blowups. The actual original-base geometry and volume equations are retained.

\end{proposition}

\begin{proof}

\uses{s7-7, s2-volume}

\leanok

Use volume descent in one direction and constructed volume-preserving elementary lift maps in the other.

\end{proof}

\begin{proposition}[7.9: positive integration equivalence]

\label{s7-9}

\lean{IntegrationInACVF.Section7.proposition_7_9_positive}

\leanok

Positive-degree volume integration and actual lifting are inverse on the relevant classes.

\end{proposition}

\begin{proof}

\uses{s7-8}

\leanok

Combine existence of volume contractions with the proved lift-kernel equivalence.

\end{proof}

\begin{theorem}[7.10: corrected all-degree integration]

\label{s7-10}

\lean{IntegrationInACVF.Section7.HK.theorem_7_10}

\leanok

Corrected graded class integration is a semiring equivalence. An arbitrary VF class integrates to the homogeneous RV blowup class of an object exactly when it is the class of that object\'s actual lift. This includes arbitrary weighted degree-zero objects.

\end{theorem}

\begin{proof}

\uses{s7-9, s7-zero-vf, s7-products}

\leanok

Combine the exact weighted zero equivalence with positive integration, then prove the actual product and coproduct laws.

\end{proof}

\begin{theorem}[7.13: the actual principal measured kernel]

\label{s7-13}

\lean{IntegrationInACVF.Section7.HK.raw_volume_kernel_principal, IntegrationInACVF.Section7.HK.theorem_7_13}

\leanok

For the corrected raw RV volume ring, the kernel of projection to the blowup ring is the ideal generated by the actual basic volume difference. Consequently the actual VF class ring is isomorphic to that single-generator quotient. The coefficient in degree one may be a weighted point.

\end{theorem}

\begin{proof}

\uses{s7-10, s7-products}

\leanok

Factor the elementary relation using corrected weighted zero coefficients, prove both ideal inclusions and apply the ring quotient equivalence. No cancellation of raw point padding is used.

\end{proof}

\section{Section 7: supported functions and integral identities}


\begin{theorem}[Full corrected coefficient linearity]

\label{s7-linear}

\lean{IntegrationInACVF.Section7.HK.gradedAEVolumeHeadIntegralLinearMap, IntegrationInACVF.Section7.HK.gradedAEVolumePrefixIntegralLinearMap, IntegrationInACVF.Section7.HK.gradedAEVolumeCoordinateIntegralLinearMap}

\leanok

Raw head, prefix and arbitrary retained-coordinate integrals are linear over the full corrected graded coefficient semiring, in every function and scalar degree. Function index $n$ represents $n+1$ VF coordinates.

\end{theorem}

\begin{proof}

\uses{s7-10, s7-products}

\leanok

Prove zero-coefficient, positive-coefficient and degree-transport identities, then descend the geometric operations to the full graded function quotient.

\end{proof}

\begin{definition}[Faithful actual support]

\label{s7-support}

\lean{IntegrationInACVF.Section7.HK.rawSupportedVolumeSubmodule, IntegrationInACVF.Section7.HK.rawSupportedVolumeRangeEquiv}

\leanok

The support submodule is exactly the faithful image of actual supported graded presentations. Coordinate and Jacobian operators preserve that support; formal support claims are tied to actual presentations.

\end{definition}

\begin{proof}

\uses{s7-linear}

\leanok



\end{proof}

\begin{proposition}[7.11: raw supported Fubini]

\label{s7-11}

\lean{IntegrationInACVF.Section7.HK.proposition_7_11}

\leanok

Coordinate integration followed by total integration equals the total integral as linear maps on all actual supported graded functions, over the full corrected coefficient semiring. Retained coordinate projections are nonempty; a separate map supplies total integration.

\end{proposition}

\begin{proof}

\uses{s7-linear, s7-support}

\leanok

Compare the forgotten supported operator with the graded geometric coordinate integral and apply total Fubini.

\end{proof}

\begin{proposition}[7.12: unpadded raw change of variables]

\label{s7-12}

\lean{IntegrationInACVF.Section7.HK.proposition_7_12}

\leanok

For an actual definable bijection, the supported Jacobian transform preserves the raw total integral as an equality of full-coefficient linear maps, in every function degree.

\end{proposition}

\begin{proof}

\uses{s7-linear, s7-support}

\leanok

Derive the supported formula from the actual geometric Jacobian map and its coefficient laws. The formula has no raw point padding to cancel.

\end{proof}

\begin{definition}[Balanced unnormalized ring functions]

\label{s7-ring}

\lean{IntegrationInACVF.Section7.HK.RawRingSupportedVolumeFunction}

\leanok

Extend the actual supported semimodule by balanced scalar extension to the corrected Grothendieck coefficient ring. Geometric coordinate and Jacobian operators extend linearly to every formal difference.

\end{definition}

\begin{proof}

\uses{s7-13, s7-support}

\leanok



\end{proof}

\begin{proposition}[7.11 over the Grothendieck ring]

\label{s7-11-ring}

\lean{IntegrationInACVF.Section7.HK.proposition_7_11_ring}

\leanok

Raw supported Fubini holds on every ring-valued function and formal difference, with the full corrected coefficient ring.

\end{proposition}

\begin{proof}

\uses{s7-ring, s7-11}

\leanok

Extend the actual linear-map identity through balanced scalar extension.

\end{proof}

\begin{proposition}[7.12 over the Grothendieck ring]

\label{s7-12-ring}

\lean{IntegrationInACVF.Section7.HK.proposition_7_12_ring}

\leanok

Raw supported change of variables holds on every ring-valued function and formal difference, without point padding or a positive function-degree restriction.

\end{proposition}

\begin{proof}

\uses{s7-ring, s7-12}

\leanok

Extend the actual supported Jacobian map and its integral identity through balanced scalar extension.

\end{proof}

\begin{lemma}[The corrected normalized commutative ring]

\label{s7-normalized}

\lean{IntegrationInACVF.Section7.HK.normalized_volume_ring_mul_comm}

\leanok

The corrected normalized coefficient ring is commutative. Normalization is performed in that ring and does not assert cancellability of a raw point class.

\end{lemma}

\begin{proof}

\uses{s7-13}

\leanok

Use the normalized basic-image commutation law and represent arbitrary elements by basic differences.

\end{proof}

\begin{theorem}[7.14: normalized supported Fubini]

\label{s7-14}

\lean{IntegrationInACVF.Section7.HK.theorem_7_14}

\leanok

For any two nonempty retained-coordinate sets of a positive-dimensional VF domain, coordinate integration followed by total integration gives the same normalized integral for every normalized supported function.

\end{theorem}

\begin{proof}

\uses{s7-11, s7-normalized}

\leanok

Both sides equal the separately defined total integral; normalize the geometric supported identity.

\end{proof}

\begin{theorem}[7.15: normalized change of variables]

\label{s7-15}

\lean{IntegrationInACVF.Section7.HK.theorem_7_15}

\leanok

The normalized integral on an actual supported domain equals the normalized integral of its Jacobian transform under any actual definable bijection, in the displayed positive-VF-dimensional function interface.

\end{theorem}

\begin{proof}

\uses{s7-12, s7-normalized}

\leanok

Descend and normalize the actual supported change-of-variables identity.

\end{proof}
